% Provenance: copied from the design workspace CambridgePQ/recursion.fast, `report/sections/C-reproducibility.tex`, on 2026-10-19. Read-only copy; the source is authoritative.
\section{Reproducibility}\label{sec:reproducibility}

Every number in this report traces to a file in the study workspace \path{recursion.fast/}. The
workspace was built on one 32-core, 62\,GB Linux machine. This appendix lists the artefacts, the
pinned revisions and tools, and the commands that regenerate each load-bearing result. Times are wall
clock on that machine, which was shared with other builds while the study ran, so they are upper
bounds. The current results are those of contract v0.3 with the round-3 fixes and of the baseline
certified against it; the v0.2 commands are kept as history.

\subsection{Workspace layout}

\begin{center}\small
\begin{tabularx}{\linewidth}{@{}>{\raggedright\arraybackslash}p{4.4cm}X@{}}
\toprule
Path & Contents \\
\midrule
\path{PLAN.md} & critique of the brief, work plan, task ledger, user inputs (U1--U3) \\
\path{notes/R1}--\path{R8} & research notes: ots.golf, sig.golf, Proximity Prize, ArkLib, leanerVM, RISC-V
  semantics, literature (with the bibliography), binary fields in Lean \\
\path{design/D1}--\path{D4} & design proposals: security and recursion; cost and incentives; binary fields;
  cost realism and parallelism \\
\path{design/DESIGN.md} & the decisions in force (v0.2 body, \S15 v0.3 amendments and round-3 fixes) \\
\path{reviews/A1}--\path{A8} & adversarial reviews: incentives, definitions, feasibility, compiled
  v0.2 contract, cryptography of the report, compiled v0.3 contract and harness, formal methods and
  editing of the report \\
\path{probes/*/} & every probe with its commands and logs; \path{probes/a6-probes/} holds A6's
  pipeline, pricing, render and hart probes \\
\path{contract-v03/} & \textbf{current}: trusted contract v0.3, library, bridge, controls, harness
  (\path{scripts/}), regressions (\path{tests/pipeline/}, \path{probes/a6/}, \path{probes/a7/}), logs
  (\path{logs/v03/}, \path{logs/orch-ci.log}), change notices \path{NOTICE-P2d.md}, \path{NOTICE-P2e.md} \\
\path{submissions/baseline-v03/} & \textbf{current}: the certified baseline (T1 and \lean{Latency1}),
  its emulator (\path{emu/}) and logs \\
\path{contract/}, \path{submissions/baseline/} & historical: contract v0.2 and the baseline certified
  against it; \path{contract/} also holds the shared Lake packages and the comparator build \\
\path{report/} & this report; \path{report/scripts/regen-appendix-a.py} regenerates and checks the
  listings of \cref{sec:appendix-contract} \\
\bottomrule
\end{tabularx}
\end{center}

\subsection{Pinned revisions and tools}

\begin{center}\small
\begin{tabularx}{\linewidth}{@{}>{\raggedright\arraybackslash}p{5.2cm}>{\raggedright\arraybackslash}X@{}}
\toprule
Component & Revision \\
\midrule
Lean toolchain & \texttt{leanprover/lean4:v4.34.1} \\
Mathlib & \texttt{v4.34.1} (\texttt{d13f23b}) \\
ArkLib (bridge and library only) & \texttt{7653a90} (leanerVM's pin; HEAD \texttt{db7005f} probed) \\
CompPoly / VCVio / PolyFun (library only) & \texttt{572f997} / \texttt{a4232d0} / \texttt{41d3b21} \\
riscv-zkvm (vendored \texttt{Rv64}) & \texttt{6810fc1}, byte-identical, \path{scripts/check-vendor.sh} \\
comparator / lean4export & \texttt{d03acab} / \texttt{076e8e5}, rebuilt on v4.34.1 (\path{contract/tools/}) \\
nanoda (second kernel) & \texttt{nanoda\_lib} v0.4.19 (\texttt{3a24072}), \path{contract-v03/tools/ext/nanoda_lib} \\
host tools (harness, baseline) & Python $\ge$3.11 (3.14.4 used), bubblewrap 0.11.1, \texttt{qemu-riscv64} 10.2.1, GNU binutils for \texttt{riscv64-unknown-elf} 2.45.50 \\
ots.golf-dev / sig.golf-dev / proximity-prize & \texttt{8b140a9} / \texttt{eaa0c7f} / \texttt{ed2b68c} \\
leanVM (Rust) & \texttt{a386121f} \\
\bottomrule
\end{tabularx}
\end{center}
The comparator, \texttt{lean4export} and nanoda binaries are not covered by the \code{.olean} hash
manifest (G-7).

\subsection{Commands for the current results (v0.3)}

\path{contract-v03/} builds only next to \path{contract/}, because its \path{lakefile.toml} sets
\texttt{packagesDir} to \path{../contract/.lake/packages}. A clean dependency setup removes that line
and the \texttt{packagesDir} key of the Lake manifest, then runs \texttt{lake exe cache get} and
\texttt{lake build}; the comparable closure measured by probe P1 took 7.8
minutes from an empty \path{.lake} (\path{probes/p1-clean-check}).

\begin{Verbatim}[fontsize=\scriptsize,frame=single,framerule=0.3pt,rulecolor=\color{black!30},framesep=4pt]
cd contract-v03
# Full CI, 10 steps (source checks, organiser package, manifest, axiom audit, score probes,
# renderer, toy through the harness, A4 and A7-M2 gate probes, pipeline regressions, manifest)
scripts/ci.sh [RISCV_ZKVM_CHECKOUT]   # round 3: CI OK, 614 s (logs/v03/ci-round3.log), independent
                                      # re-run 624 s (logs/orch-ci.log); round 4: CI OK 685 s
                                      # (logs/v03/ci-round4.log), independent re-run 687 s (logs/orch-ci-r4.log)
scripts/org-build.sh .lake/org        # read-only organiser package used by every harness stage
python3 -I scripts/olean-manifest.py check --org .lake/org scripts/olean-manifest.sha256  # 35,093 files
(cd .lake/org && lake env .lake/build/bin/axiomaudit)                  # logs/v03/axioms.log
python3 scripts/loc.py RecursionFast/Contract.lean RecursionFast/Contract/*.lean   # 796 / 471 / 1,536
G=.lake/org/.lake/build/bin/gate
tests/pipeline/cmpproj.sh; tests/pipeline/gateproj.sh $G   # A6-G0, G1/G2 regressions (CI step 9)
tests/pipeline/bindproj.sh $G      # A7-M3: gate-tree build must equal the comparator's (round 4)

# Production harness on one submission (stage 0 manifest; 1 comparator first, fresh tree;
# 2 gate, second fresh tree, build byte-identical to stage 1; 3 optional --badge CMD as its own process)
scripts/harness.sh --org .lake/org --manifest scripts/olean-manifest.sha256 \
  --gate RecursionFast.Challenge.certified --computable RecursionFast.Challenge.submission \
  CHALLENGE.lean SOLUTION_DIR config.json        # prints HARNESS: ACCEPT or HARNESS: REJECT (stage)

# Probes against the contract's own functions
lake env lean probes/a6/Pricing.lean                     # A6-P1 leaf pricing (also CI step 5)
lake env lean ../probes/a6-probes/EpsComputable.lean     # A6-N3: the v0.3 limitation (fixed in round 4;
                                                         # regression probes/a7/m2/M2Err.lean, CI step 8)

# Certified baseline (T1 and Latency1), from the workspace root; needs riscv64-unknown-elf-as
submissions/baseline-v03/scripts/check.sh          # CHECK OK: 815 s; independent re-run 13 min 35 s
submissions/baseline-v03/emu/measure.sh 16 20 24   # class runs at the nonzero cap (k=24: ~10 GB, ~40 min)

# Appendix A against the contract source (from report/)
python3 -I scripts/regen-appendix-a.py ../contract-v03 sections/A-lean-contract.tex --check
\end{Verbatim}

\subsection{Historical commands (v0.2)}

\begin{Verbatim}[fontsize=\scriptsize,frame=single,framerule=0.3pt,rulecolor=\color{black!30},framesep=4pt]
# Trusted contract v0.2: build, vendor check, hygiene, controls, axiom audit, comparator
cd contract && lake build && bash scripts/ci.sh             # ~72 s warm
# Baseline certified against v0.2 (image, hygiene, build, axioms, challenge, comparator, emulator diff)
bash submissions/baseline/scripts/check.sh                 # "CHECK OK", 2 min 15 s warm
# Round-1 counterexamples (from a leanerVM checkout) and round-2 probes (A4, compiled v0.2 contract)
lake env lean <recursion.fast>/probes/a2-probes/A2Probes.lean
cd contract && lake env lean ../probes/a4-probes/<Probe>.lean
# Cost-model probes (nine models; field kernels; trajectory)
bash probes/d4-probes/run.sh
python3 probes/d3-probes/model/trajectory_bin.py
\end{Verbatim}

\subsection{What was not re-run}

The following were measured once by the probe that produced them:
\begin{itemize}[nosep]
\item the sig.golf candidate build and replay (\path{probes/sig-golf-build});
\item the ots.golf record check (R1);
\item the Proximity Prize builds (\path{probes/proximity-build});
\item review A7's scratch probes (\lean{ε} computability, an IR-overwrite attempt), which were not kept
  in the workspace; \path{probes/a6-probes/EpsComputable.lean} reproduces the first, and the
  round-4 regressions \path{contract-v03/probes/a7/} cover A7-M2 and M3.
\end{itemize}
The literature claims carry the verification tags of \path{notes/R7-literature.md} and the source
checks of review A5.
